% 第1回　ガイダンス
\session{1}{1}{9月30日（水）}{ガイダンス}{}{}

\goals{
  \item 授業の二部構成と評価の仕組みを説明できる
  \item この授業でのAI利用のルールを説明できる
  \item ハルシネーションがなぜ起きるかを理解し、出力を確かめる方法を挙げられる
}

\work{確認}{授業の全体像と評価}
\begin{wstab}{colspec={Q[l,wd=40mm,m] X[l,m]},row{1-Z}={ht=9mm}}
  \hc{前半（第1〜5回）} & \\
  \hc{後半（第6〜13回）} & \\
  \hc{成績評価の内訳} & 提出物 \underline{\hspace{12mm}}\,\%　／　発表 \underline{\hspace{12mm}}\,\% \\
  \hc{主な提出物} & \\
\end{wstab}

\work{確認}{この授業でのAI利用のルール（自分の言葉で3つ）}
\writelines{3}

\work[12mm]{準備}{Copilot の動作確認}
\begin{checks}
  \item 大学のアカウントで Copilot にサインインできた
  \item 簡単な質問をして、回答が返ってくることを確かめた
\end{checks}

\work[70mm]{ミニワーク}{AIに嘘をつかせる}
\begin{promptbox}[ハルシネーションを体験する]
名古屋市立大学経済学部の歴史について、参考文献を3つ挙げて300字で説明してください。
\end{promptbox}
\instr{(1) AIが挙げた参考文献が実在するか、図書館の蔵書検索や CiNii Research で確かめる。}
\begin{wstabh}{colspec={X[3,l,m] Q[c,wd=14mm,m] X[2,l,m]},row{2-Z}={ht=10mm}}
  AIが挙げた文献（著者・題名・年） & 実在 ○× & 確かめた場所・わかったこと \\
  & & \\ & & \\ & & \\
\end{wstabh}

\instr{(2) 自分が詳しく知っている話題（出身地、部活、趣味など）について聞き、内容の正誤を判定する。}
\begin{wstabh}{colspec={X[2,l,m] X[3,l,m] Q[c,wd=14mm,m] X[2,l,m]},row{2-Z}={ht=12mm}}
  尋ねたこと & AIの答え（要点） & 正誤 ○△× & どこが違ったか \\
  & & & \\ & & & \\ & & & \\
\end{wstabh}

\work{まとめ}{気づいたこと・「出力を使う前のチェック」}
\instr{共有で出た意見も含め、AIの出力を使う前に確かめるべきことを自分の言葉で整理する。}
\writelines{5}

\begin{nexttime}
  \item 関心のあるテーマを3つ書き出す（第2回のブレストで使う）。「なぜ気になるのか」を一言添える。
\end{nexttime}
\begin{wstabh}{colspec={Q[c,wd=8mm,m] X[2,l,m] X[3,l,m]},row{2-Z}={ht=11mm}}
  & 関心のあるテーマ & なぜ気になるのか \\
  1 & & \\ 2 & & \\ 3 & & \\
\end{wstabh}
